\section*{Chapter \ref{ch:s2:convertible-arbitrage} --- Convertible Arbitrage}

\begin{solution}{exo:s2:convertible-arbitrage:1}
$100/2.5 = 40$ a share; at 60 the conversion value is $2.5 \times 60 = 150$.
\end{solution}

\begin{solution}{exo:s2:convertible-arbitrage:2}
Its bonds are three times its capital, so $3 \times 2.7 = 8.1\%$ of capital, close to the $-7.2\%$ the funds reported for January to May 2005.
\end{solution}

\begin{solution}{exo:s2:convertible-arbitrage:3}
$1.27 \times 15/81.2 = 0.23$: a 1\% fall in the share takes 0.23\% off the bond.
\end{solution}

\begin{solution}{exo:s2:convertible-arbitrage:4}
The straight bond is discounted at the spread for the issue's hazard at a share price of 40. At 15 the equity-to-credit link has raised the hazard,
so the bond is worth less than a bond at the old spread; the conversion option is too far out of the money to make up for it.
\end{solution}

\begin{solution}{exo:s2:convertible-arbitrage:5}
Convexity: as the share fell, the bond's delta fell, so it lost less per unit of fall than the short hedge, sized at the higher delta, gained. The book
is long gamma, and the crash's large moves paid 3.7 points net.
\end{solution}

\begin{solution}{exo:s2:convertible-arbitrage:6}
Cheapness is the gap between market price and model value, driven by who must sell and who can buy. The share, the credit default swap and the
option market do not carry it, so no traded hedge offsets it; only the ability to hold until it closes does.
\end{solution}

\begin{solution}{exo:s2:convertible-arbitrage:7}
Over five years the credit hedge costs 21.8\% of capital in premiums and gives back the credit rebound; it lowers the carry from 9.8\% to 6.0\% a
year and the five-year return from 51.9\% to 30.1\%. It buys the episode's 6.9 points of credit loss, halving the episode from $-11.6\%$ to
$-5.8\%$, and leaves the 12.6 points of cheapening untouched.
\end{solution}

\begin{solution}{exo:s2:convertible-arbitrage:8}
Delta and credit hedges leave cheapening, borrow and financing risk: in the chapter's episode cheapening alone cost 12.6\% of capital at three to one,
so at ten to one it would cost about 42\%, before any margin call. In 2005 and 2008 the risk that mattered was exactly the one the hedges left.
\end{solution}

\begin{solution}{pb:s2:convertible-arbitrage:1}
\begin{enumerate}
\item A straight bond, an option on the shares, less the issuer's call, all exposed to default.
\item 117.7 against a conversion value of 100, a delta of 1.80 shares.
\item A convertible far below its conversion price that trades on credit; at 15 it is worth 81.2, below the straight bond's 87.34, with a delta of 1.27
and an elasticity of 0.23.
\item The hedge's delta includes the credit effect of the share price; a hedge that ignores it is too small when the share falls.
\item As defined in the chapter: a short stock position equal to the delta; plus credit protection sized to the credit sensitivity.
\item A buy-and-hedge strategy explains much of the funds' returns; extreme events and convertible supply matter.
\item Redemptions of more than 20\% in a quarter, forced sales, a 2.7\% median discount in May, $-7.2\%$ fund returns.
\item The leverage offered by prime brokers fell suddenly as they faced failure.
\item Thirty bonds at issue, short their deltas daily, levered three to one; P\&L in gamma, credit, cheapness, coupons and borrow.
\item 9.8\% a year delta-hedged; 6.0\% credit-hedged.
\item $-11.6\%$: cheapness $-12.6$, credit $-6.9$, gamma $+3.7$, coupon $+5.3$, borrow $-1.1$.
\item The bonds recovered their cheapness and the hazard fell back: $+22.0\%$ in six months.
\item \textbf{Named result.} The delta-hedged book carries 9.8\% a year on capital and loses 11.6\% in the planted three-month episode, 6.9 points from
credit and 12.6 from cheapening; hedging credit halves the loss to 5.8\% at the cost of carry (6.0\% a year).
\item It costs 21.8\% of capital over five years and buys the 6.9 points of credit loss.
\item Cheapening is multiplied by leverage and can force selling at the worst prices.
\item They faced redemptions (2005) or lost leverage (2008): capital, not prices, decided.
\item By the loss in a forced-selling episode at the chosen leverage, against the capital that could be lost without a forced unwind.
\item Model values with realistic inputs, episodes of cheapening, borrow and financing costs that change.
\item The new-issue convertible.
\item For funding issuers quickly and holding illiquid bonds through the times when others must sell.
\end{enumerate}
\end{solution}

\begin{solution}{iq:s2:convertible-arbitrage:1}
Buy the convertible, short the delta in shares, possibly buy credit protection; earn the coupon, the short rebate less borrow, the gamma from
rehedging and the convergence of a cheap bond to its value; risks are cheapening, credit, borrow and financing.
\end{solution}

\begin{solution}{iq:s2:convertible-arbitrage:2}
When the share falls the hazard rises, so the bond falls more than an option model with constant credit says: the delta is larger at low share
prices, and at very low prices most of the delta is credit.
\end{solution}

\begin{solution}{iq:s2:convertible-arbitrage:3}
Cheapening, borrow recall, financing withdrawal, gamma in gaps, issuer calls and takeovers, model error in the credit input; stress with a 2005-style
cheapening, a doubled hazard and a borrow recall at once, at the book's leverage.
\end{solution}

\begin{solution}{iq:s2:convertible-arbitrage:4}
Find borrow elsewhere, hedge with options or a correlated instrument, or reduce the position; accept a less precise hedge rather than an unhedged
bond, and record the cost for the book's borrow risk.
\end{solution}

\begin{solution}{iq:s2:convertible-arbitrage:5}
Each day: prices of bonds, shares, credit and vols; model values and deltas from the pricer (tabulated for speed); hedge adjustments net across
issuers; borrow availability; cheapness monitoring and attribution; checks on corporate actions, calls and conversions.
\end{solution}

\begin{solution}{iq:s2:convertible-arbitrage:6}
As for any option, the hedged position's value changes by $\Theta\,\mathrm{d}t + \tfrac12\Gamma(\mathrm{d}S)^2$, with the bond's coupon and the short's
rebate less borrow added as carry. At default the share jumps to near zero: the short stock gains its value and the bond falls to its recovery; with
an equity-to-credit model the delta hedge has been sized for part of that jump, and credit protection covers the rest.
\end{solution}
